% =============================================================================
% Appendix A: Blinding Protocol
% =============================================================================

\section{Blinding Protocol}
\label{sec:blinding}

% AGENT: Document the specific blinding variable and signal region definition
% from analysis strategy. Read blinding configuration from
% gad/config/analysis_config.yaml and strategy document.

This appendix documents the blinding protocol followed during this analysis,
including the staged unblinding sequence, pre-unblinding checklist results,
signal injection tests, and blinding integrity verification.

\subsection{Staged Unblinding Sequence}
\label{sec:blinding:stages}

The analysis follows a staged unblinding protocol with four stages, designed
to minimize confirmation bias and ensure the robustness of the analysis
before examining signal-region data.

\begin{enumerate}
  \item \textbf{BLINDED:} Full blinding. Signal region data is never accessed.
        All analysis development uses sideband data and simulation only.
        Plots in the signal region show MC-only expectations.

  \item \textbf{SIDEBANDS:} Sideband unblinding. Data in control and validation
        regions is compared with predictions. No signal region data is
        examined. Background estimation methods are validated.

  \item \textbf{PARTIAL\_10PCT:} Partial unblinding with 10\% of signal region
        data. A randomly selected 10\% subset of SR data is examined to
        check for gross mismodelling. The subset selection is deterministic
        (based on event number modulo) and does not bias the final result.

  \item \textbf{UNBLINDED:} Full unblinding. The complete signal region dataset
        is used for the final statistical analysis. This stage is only
        entered after all pre-unblinding checks have passed.
\end{enumerate}

% AGENT: Document the specific blinding variable and signal region definition
% from the analysis strategy. Include which branches/variables define the SR
% boundary and how the blinding mechanism prevents data access.

The blinding mechanism is enforced at the data loading level through the
\texttt{BlindingManager} class, which raises a \texttt{BlindingViolationError}
if signal region data is accessed in an unauthorized blinding state.

\subsection{Pre-Unblinding Checklist}
\label{sec:blinding:checklist}

% AGENT: Read unblinding_checklist from STATE.md frontmatter.
% Format as table with Item, Status, Evidence columns.
% Each checklist item must be marked PASS/FAIL with supporting evidence.

Before proceeding to full unblinding, the following checklist items must
all be satisfied. The status of each item is recorded in the analysis
state file and verified by the cross-checker agent.

\begin{table}[htbp]
  \centering
  \caption{Pre-unblinding checklist status. All items must pass before
           proceeding to full unblinding.}
  \label{tab:checklist}
  \begin{tabular}{p{5cm}lp{6cm}}
    \toprule
    Item & Status & Evidence \\
    \midrule
    % AGENT: Read checklist items from STATE.md unblinding_checklist field.
    % Format each item as a row. Example items:
    Background model closure in VRs
      & [PASS/FAIL]
      & [chi2/ndf values, closure test results] \\
    \addlinespace
    Systematic uncertainties evaluated
      & [PASS/FAIL]
      & [Number of NPs, coverage tests] \\
    \addlinespace
    Expected sensitivity computed
      & [PASS/FAIL]
      & [Expected limit value, Asimov fit status] \\
    \addlinespace
    Fit diagnostics satisfactory
      & [PASS/FAIL]
      & [NP pulls within bounds, no large constraints] \\
    \addlinespace
    Signal injection test passed
      & [PASS/FAIL]
      & [Pull distribution mean, width] \\
    \addlinespace
    Cross-check analysis consistent
      & [PASS/FAIL]
      & [Limit comparison, delta/tolerance] \\
    \addlinespace
    Analysis note pre-unblinding draft complete
      & [PASS/FAIL]
      & [Sections 1--8 reviewed] \\
    \bottomrule
  \end{tabular}
\end{table}

\subsection{Signal Injection Tests}
\label{sec:blinding:injection}

% AGENT: Read injection test results from Wave 5 cross-checker report.
% Format as table with mu_injected, mu_hat, sigma_fit, pull, pass/fail.

Signal injection tests are performed to validate that the statistical model
can correctly recover an injected signal. Asimov datasets are generated with
various injected signal strengths, and the fitted signal strength
$\hat{\mu}$ is compared with the injected value.

\begin{table}[htbp]
  \centering
  \caption{Signal injection test results. The pull is defined as
           $(\hat{\mu} - \mu_\text{inj}) / \sigma_{\hat{\mu}}$.
           All pulls should be within $\pm 2$.}
  \label{tab:injection}
  \begin{tabular}{rrrrr}
    \toprule
    $\mu_\text{inj}$ & $\hat{\mu}$ & $\sigma_{\hat{\mu}}$
      & Pull & Status \\
    \midrule
    % AGENT: Fill from signal injection test results
    0.0 & -- & -- & -- & [PASS/FAIL] \\
    0.5 & -- & -- & -- & [PASS/FAIL] \\
    1.0 & -- & -- & -- & [PASS/FAIL] \\
    2.0 & -- & -- & -- & [PASS/FAIL] \\
    \bottomrule
  \end{tabular}
\end{table}

\subsection{Blinding Integrity Verification}
\label{sec:blinding:integrity}

% AGENT: Read integrity check results from Wave 5 cross-checker report.
% Document checks performed on SR plots, wave outputs, and git history.

The blinding integrity is verified through the following automated checks:

\begin{itemize}
  \item \textbf{Signal region plot audit:} All plots in the analysis output
        directories are scanned to verify that no SR data distributions
        are present before the appropriate unblinding stage.
        % AGENT: Fill result: number of plots checked, violations found.

  \item \textbf{Wave output audit:} All wave output files are checked to
        ensure no SR data values appear in reports or logs.
        % AGENT: Fill result: number of files checked, violations found.

  \item \textbf{Git history audit:} The git commit history is scanned to
        verify that no commits contain SR data in committed files.
        % AGENT: Fill result: number of commits checked, violations found.

  \item \textbf{Blinding state transitions:} The sequence of blinding state
        transitions is verified against the expected protocol.
        % AGENT: Fill from STATE.md blinding state transition log.
\end{itemize}

\subsection{Re-Blinding Events}
\label{sec:blinding:reblinding}

% AGENT: Read reblind_count and post_unblinding_problems from STATE.md.
% If reblind_count > 0, describe the event: when it occurred, why
% re-blinding was triggered, what was fixed, and the subsequent
% re-unblinding procedure.

% AGENT: If reblind_count == 0, keep the following statement:
No re-blinding events occurred during this analysis. The blinding protocol
was followed without interruption from initial blinding through to final
unblinding.

% AGENT: If reblind_count > 0, replace the above with a description of
% each re-blinding event in the following format:
% \begin{enumerate}
%   \item \textbf{Re-blinding event [N]:}
%     \begin{itemize}
%       \item \textbf{Date:} [date from STATE.md]
%       \item \textbf{Trigger:} Category B problem -- [description]
%       \item \textbf{Resolution:} [what was fixed]
%       \item \textbf{Re-unblinding:} [when/how re-unblinding was performed]
%     \end{itemize}
% \end{enumerate}

\subsection{Summary}
\label{sec:blinding:summary}

% AGENT: Provide a brief summary of the blinding protocol outcome.
% State the final blinding status, number of re-blinding events,
% and overall assessment of blinding integrity.

The blinding protocol was followed throughout this analysis. The analysis
progressed through stages BLINDED $\to$ SIDEBANDS $\to$ PARTIAL\_10PCT
$\to$ UNBLINDED with all pre-unblinding checks passing before each
transition. [N] re-blinding events occurred. The blinding integrity
verification found no violations.
